Un condensateur à air de \qty{\v{C0}}{nF} est chargé par une pile de \qty{\v{V0}}{V}, puis débranché de la pile. On insère ensuite entre ses armatures un diélectrique de constante $\kappa = \v{kap}$ qui remplit tout l'espace.

\begin{enumerate}[label=\alph*)]
    \item Déterminez la charge et l'énergie du condensateur avant l'insertion du diélectrique.
    \item Déterminez la nouvelle capacité et la nouvelle différence de potentiel.
    \item Déterminez la nouvelle énergie emmagasinée. Le diélectrique est-il attiré ou repoussé lors de l'insertion ? Justifiez.
\end{enumerate}
